% Options for packages loaded elsewhere
\PassOptionsToPackage{unicode}{hyperref}
\PassOptionsToPackage{hyphens}{url}
\documentclass[
  ignorenonframetext,
]{beamer}
\newif\ifbibliography
\usepackage{pgfpages}
\setbeamertemplate{caption}[numbered]
\setbeamertemplate{caption label separator}{: }
\setbeamercolor{caption name}{fg=normal text.fg}
\beamertemplatenavigationsymbolsempty
% remove section numbering
\setbeamertemplate{part page}{
  \centering
  \begin{beamercolorbox}[sep=16pt,center]{part title}
    \usebeamerfont{part title}\insertpart\par
  \end{beamercolorbox}
}
\setbeamertemplate{section page}{
  \centering
  \begin{beamercolorbox}[sep=12pt,center]{section title}
    \usebeamerfont{section title}\insertsection\par
  \end{beamercolorbox}
}
\setbeamertemplate{subsection page}{
  \centering
  \begin{beamercolorbox}[sep=8pt,center]{subsection title}
    \usebeamerfont{subsection title}\insertsubsection\par
  \end{beamercolorbox}
}
% Prevent slide breaks in the middle of a paragraph
\widowpenalties 1 10000
\raggedbottom
\AtBeginPart{
  \frame{\partpage}
}
\AtBeginSection{
  \ifbibliography
  \else
    \frame{\sectionpage}
  \fi
}
\AtBeginSubsection{
  \frame{\subsectionpage}
}
\usepackage{iftex}
\ifPDFTeX
  \usepackage[T1]{fontenc}
  \usepackage[utf8]{inputenc}
  \usepackage{textcomp} % provide euro and other symbols
\else % if luatex or xetex
  \usepackage{unicode-math} % this also loads fontspec
  \defaultfontfeatures{Scale=MatchLowercase}
  \defaultfontfeatures[\rmfamily]{Ligatures=TeX,Scale=1}
\fi
\usepackage{lmodern}
\ifPDFTeX\else
  % xetex/luatex font selection
\fi
% Use upquote if available, for straight quotes in verbatim environments
\IfFileExists{upquote.sty}{\usepackage{upquote}}{}
\IfFileExists{microtype.sty}{% use microtype if available
  \usepackage[]{microtype}
  \UseMicrotypeSet[protrusion]{basicmath} % disable protrusion for tt fonts
}{}
\makeatletter
\@ifundefined{KOMAClassName}{% if non-KOMA class
  \IfFileExists{parskip.sty}{%
    \usepackage{parskip}
  }{% else
    \setlength{\parindent}{0pt}
    \setlength{\parskip}{6pt plus 2pt minus 1pt}}
}{% if KOMA class
  \KOMAoptions{parskip=half}}
\makeatother
\setlength{\emergencystretch}{3em} % prevent overfull lines
\providecommand{\tightlist}{%
  \setlength{\itemsep}{0pt}\setlength{\parskip}{0pt}}
% Slide layout defaults for warning-clean scientific decks.
\usepackage{etoolbox}
\IfFileExists{xurl.sty}{\usepackage{xurl}}{}
\IfFileExists{seqsplit.sty}{\usepackage{seqsplit}}{\newcommand{\seqsplit}[1]{#1}}
\protected\def\breakseq#1{\seqsplit{#1}}
\protected\def\breaktt#1{\begingroup\ttfamily\seqsplit{#1}\endgroup}
\setlength{\emergencystretch}{6em}
\tolerance=5000
\hbadness=10000
\hfuzz=1pt
\setlength{\tabcolsep}{2pt}
\AtBeginEnvironment{longtable}{\tiny\renewcommand{\arraystretch}{0.86}\setlength{\tabcolsep}{1pt}}
\AtBeginEnvironment{tabular}{\tiny\renewcommand{\arraystretch}{0.86}\setlength{\tabcolsep}{1pt}}

% Natbib and cross-reference fallbacks — slides don't load natbib
% or cleveref, but combined-PDF manuscript prose may emit these
% commands. The fallback renders citations as a bracketed key list
% and cross-references as detokenized labels so slides don't fail on
% undefined control sequences or raw underscores. \providecommand is
% a no-op if packages load later.
\providecommand{\citep}[1]{[#1]}
\providecommand{\citet}[1]{#1}
\providecommand{\citealp}[1]{#1}
\providecommand{\citeauthor}[1]{#1}
\providecommand{\citeyear}[1]{#1}
\providecommand{\cref}[1]{\texttt{\detokenize{#1}}}
\providecommand{\Cref}[1]{\texttt{\detokenize{#1}}}

\newtheorem{proposition}{Proposition}
\newtheorem{hypothesis}{Hypothesis}
\usepackage{bookmark}
\IfFileExists{xurl.sty}{\usepackage{xurl}}{} % add URL line breaks if available
\urlstyle{same}
% fallback for those not using the hyperref driver hyperxmp:
\makeatletter
\@ifundefined{xmpquote}{\newcommand{\xmpquote}[1]{#1}}{}
\makeatother
\hypersetup{
  hidelinks,
  pdfcreator={LaTeX via pandoc}}

\author{\texorpdfstring{}{}}
\date{}

\begin{document}

\begin{frame}[fragile,allowframebreaks]{Stage 2: Bibliometric Analysis
\label{sec:methods_bibliometrics}}
\protect\phantomsection\label{stage-2-bibliometric-analysis}
Stage 2 performs four complementary analyses on the deduplicated corpus.
All analyses are deterministic given fixed random seeds and operate on
the same \texttt{corpus.jsonl} input.

\begin{block}{Subfield Classification}
\protect\phantomsection\label{subfield-classification}
Each paper is classified into one of eight categories organized across
three domains: \textbf{A -- Core Theory} (A1: quantitative and formal
mathematical theory; A2: qualitative philosophy and general FEP theory),
\textbf{B -- Tools \& Translation} (algorithms, scaling, and software
development), and \textbf{C -- Application Domains} (C1: neuroscience,
C2: robotics, C3: language processing, C4: computational psychiatry, C5:
biology and morphogenesis). Classification uses word-boundary-aware
keyword matching against curated lists (74+ mathematical indicators, 25+
philosophy terms, 24+ tools terms, and 14--20 terms per application
domain---totaling over 200 keywords across 8 categories, all documented
in \texttt{config.yaml}) applied to titles and abstracts. A priority
system ensures that specific application domains (C1--C5, priority 1)
take precedence over tools (B, priority 2), formal theory (A1, priority
3), and the broad qualitative philosophy catch-all (A2, priority 4).
Within a priority tier, the domain with the most keyword matches wins.
A1's keyword set includes mathematical indicators such as
\emph{theorem}, \emph{proof}, \emph{convergence}, \emph{posterior},
\emph{equation}, and \emph{Fokker--Planck}, ensuring that papers with
mathematical content are classified as formal theory rather than
defaulting to the philosophy category.
\end{block}

\begin{block}{Temporal Metrics and Growth-Rate Estimation}
\protect\phantomsection\label{temporal-metrics-and-growth-rate-estimation}
We compute temporal publication metrics including year-by-year counts
with gap-filling, cumulative totals, 3-year smoothed moving averages,
and peak year identification. Field dynamics are estimated via two
complementary metrics. The \textbf{mean year-over-year growth rate}
\(\bar{g}\) is the arithmetic mean of annual growth rates for years with
non-zero prior-year publications. The \textbf{doubling time}
\(t_d = \ln 2 / \ln(1 + \bar{g})\). The \textbf{compound annual growth
rate} (CAGR) captures the annualized rate across the full temporal span.
Mathematical details are provided in Appendix \ref{sec:appendix_growth}.
\end{block}

\begin{block}{Text Analytics}
\protect\phantomsection\label{text-analytics}
We construct the TF-IDF matrix using tokenization with stopword removal
and L2-normalized smoothed term-frequency inverse-document-frequency
weighting \citep{salton1975vector}, with a configurable vocabulary size
(default: 500 features). We apply non-negative matrix factorization
(NMF) to discover latent topics using multiplicative update rules
\citep{lee1999nmf}. Topic count \(k=5\) was selected via expert-driven
assessment. Mathematical details are provided in Appendix
\ref{sec:appendix_nmf}.
\end{block}

\begin{block}{Citation Network Construction}
\protect\phantomsection\label{citation-network-construction}
We construct the intra-corpus citation network as a directed graph where
nodes are papers and edges represent citation relationships resolved
within the corpus. Because identifier formats vary across databases
(arXiv IDs, DOIs, Semantic Scholar IDs), only references whose
identifiers match a corpus entry contribute edges; the resulting
resolution rate (7.4\%) represents a lower bound on the true
intra-corpus citation density. Network metrics include PageRank
centrality, HITS hub and authority scores
\citep{kleinberg1999authoritative}, degree distributions, network
density, connected components, and community structure via greedy
modularity maximization \citep{clauset2004finding}.
\end{block}
\end{frame}

\end{document}
